\begin{center}\begin{tikzpicture}[font=\small,>=stealth,line join=round]
\draw[gray!40,rounded corners=3pt] (0,0) rectangle (14.2,8.3);\node at (7.1,7.700000000000001) {例5：把已知的右半边，关于竖线配到左侧};
\begin{scope}[shift={(1.7000000000000002,1.8921052631578945)},x=1.2cm,y=0.2421052631578947cm]\draw[->,gray] (-0.5,0)--(4.5,0) node[right]{$x$};\draw[->,gray] (0,-1)--(0,18) node[above]{$y$};\draw[violet,thick,dashed,domain=0:2,samples=180] plot(\x,{pow(2,4-\x)});\draw[blue,thick,solid,domain=2:4,samples=180] plot(\x,{pow(2,\x)});\fill[orange!85!black] (0,16) circle(2.2pt);\fill[orange!85!black] (1,8) circle(2.2pt);\fill[orange!85!black] (2,4) circle(2.2pt);\fill[orange!85!black] (3,8) circle(2.2pt);\fill[orange!85!black] (4,16) circle(2.2pt);\draw[gray,dashed] (2,-.5)--(2,17);\end{scope}\node[anchor=west,align=left,text width=5.0cm] at (8.4,4.0) {蓝：$x\ge2,\ f(x)=2^x$\\[9pt]紫：$x<2,\ f(x)=2^{4-x}$\\[9pt]$0\leftrightarrow4$，同为 $16$\\[9pt]$1\leftrightarrow3$，同为 $8$\\[9pt]$4-x>2$ 才能用右式};
\node[text width=12.6cm,align=center] at (7.1,.65) {图示区间 $[0,4]$；完整函数定义在实数上，左侧指数中的自变量是 $4-x$。};
\end{tikzpicture}\end{center}
